\documentclass[11pt]{article}

\usepackage[utf8]{inputenc}
\usepackage[T1]{fontenc}
\usepackage{lmodern}
\usepackage[margin=1in]{geometry}
\usepackage{booktabs}
\usepackage{longtable}
\usepackage{amsmath}
\usepackage{amssymb}
\usepackage{xcolor}
\usepackage{listings}
\usepackage{hyperref}
\hypersetup{colorlinks=true, linkcolor=blue!55!black, urlcolor=blue!55!black,
            citecolor=blue!55!black}

\definecolor{codebg}{RGB}{245,245,245}
\lstset{
  basicstyle=\ttfamily\small,
  backgroundcolor=\color{codebg},
  breaklines=true,
  frame=single,
  rulecolor=\color{gray!40},
  columns=fullflexible,
  keepspaces=true,
  showstringspaces=false,
}

\newcommand{\code}[1]{\texttt{#1}}

\title{\textbf{Critical Notes: Data-Provenance Defects and Fixes}\\[0.4em]
       \large KG-Commit Reproducibility Package}
\author{}
\date{Recorded 2026-07-14}

\begin{document}
\maketitle

\begin{abstract}
This document records a critical data-provenance defect discovered in the
cross-project build of the KG-Commit pipeline, the investigation that isolated
its root cause, the decision taken to fix it, and the engineering logic of that
fix. It is written so that a future reader (or reviewer) can understand exactly
what went wrong, why the first \code{hadoop-mapreduce} graph was silently
incomplete, why the chosen remedy (a \emph{multi-repository} build) is the
faithful one, and what limitations remain. The same failure mode can affect
\emph{any} project whose defect labels were mined across a repository split, so
the note is intended as a durable reference, not a one-off changelog entry.
\end{abstract}

\tableofcontents

% ======================================================================
\section{Summary}
% ======================================================================
The first end-to-end build of the \code{hadoop-mapreduce} project silently
produced a knowledge graph covering only \textbf{525 of the 1{,}321 labelled
commits} (39.7\%). Every layer was affected---not only the semantic text
(CSTG) layer but also the structural layers (AST delta, CFG, DFG, PDG, SEQ),
which are the most critical to the method. The build \emph{reported success}
because the missing commits produced no error the driver checked for: their
commit SHAs simply could not be resolved in the git repository the build was
pointed at, so they were skipped.

The root cause is a \textbf{repository-split provenance mismatch}. The ApacheJIT
labels for ``hadoop-mapreduce'' were mined from commit histories that are today
distributed across \emph{two} distinct git repositories, with \textbf{zero SHA
overlap}. A single-repository build can therefore never see the whole labelled
stream. The fix is to make the build engine resolve each commit from
whichever repository contains it (a \emph{multi-repository} build), then rebuild
the project on the full 1{,}321-commit stream.

% ======================================================================
\section{How the defect surfaced}
% ======================================================================
Two independent symptoms pointed at the same underlying problem:

\begin{enumerate}
  \item \textbf{Frozen progress at 39.7\%.} The per-commit build timing logs
  under \code{logs/hadoop-mapreduce/} contained only 525 of the expected 1{,}321
  rows for \emph{every} online-growth layer (AST/CFG/DFG/PDG/SEQ). This was
  initially mistaken for a logging artifact; it was in fact the literal truth of
  what the build processed.

  \item \textbf{Missing diff text.} The per-project diff CSV
  \code{data/apachejit/projects/apache\_hadoop-mapreduce\_diff.csv} had empty
  \code{diff\_text} in 796 of 1{,}321 rows. Attempting to regenerate them from
  the cloned \code{apache/hadoop-mapreduce} repository failed: git returned
  \code{SHA <...> could not be resolved} for exactly those 796 commits.
\end{enumerate}

The number \emph{796} appearing in both symptoms ($796 + 525 = 1{,}321$) was the
clue that a single cause underlay both.

% ======================================================================
\section{Root-cause investigation}
% ======================================================================

\subsection{The two repositories}
\begin{center}
\begin{tabular}{@{}lrl@{}}
\toprule
Repository & Commits & History span \\
\midrule
\code{apache/hadoop-mapreduce} (spin-off) & 960     & 2009-05-19 \ldots{} 2011-06-12 \\
\code{apache/hadoop} (monorepo)           & 28{,}289 & 2009 \ldots{} 2019 (and ongoing) \\
\bottomrule
\end{tabular}
\end{center}

\code{apache/hadoop-mapreduce} is a \emph{short-lived spin-off} repository. During
2009--2011 the Hadoop project temporarily split its monolithic repository into
sub-project repositories (\code{hadoop-common}, \code{hadoop-hdfs},
\code{hadoop-mapreduce}). In 2011 these were folded back into the single
\code{apache/hadoop} monorepo. The re-merge \textbf{rewrote history}: the same
logical changes acquired \emph{new} commit SHAs in the monorepo, so the spin-off
SHAs and the monorepo SHAs do not coincide.

\subsection{The label partition}
Resolving each of the 1{,}321 ApacheJIT-labelled commit SHAs against both
repositories yields a clean, non-overlapping partition:

\begin{center}
\begin{tabular}{@{}lr@{}}
\toprule
Resolves in \ldots{} & Commits \\
\midrule
spin-off repository \emph{only} & 525 \\
monorepo \emph{only}            & 796 \\
both repositories               & 0 \\
neither repository              & 0 \\
\midrule
total labelled                  & 1{,}321 \\
\bottomrule
\end{tabular}
\end{center}

Consequently \emph{no single repository contains the full labelled set}. The
original build used the spin-off repository (\code{repo\_subpath:
repos/apache/hadoop-mapreduce}); it saw the 525 spin-off-only commits and
silently skipped the 796 monorepo-only ones.

\subsection{Why every layer was affected, not only CSTG}
It is important to be precise about which inputs each layer reads, because it
would be tempting (and wrong) to conclude that only the text/CSTG layer was hurt.

\begin{itemize}
  \item The \textbf{structural layers} (AST delta, CFG, DFG, PDG, SEQ) are built
  by \emph{replaying commits from the git repository}: for each commit the engine
  runs \code{git show <ref>:<path>} and \code{git diff --name-status} to recover
  the changed files and their before/after blobs
  (\code{build/build\_online\_kg.py}, \code{build/build\_subgraph\_online\_kg.py}).
  If the commit SHA does not resolve in the repository, the commit contributes
  \emph{no nodes and no edges} to any structural layer.

  \item The \textbf{CSTG layer} reads commit text from the diff CSV
  (\code{diff\_text}); those rows were also empty for the 796.
\end{itemize}

A direct query against the built graph confirmed the damage: of the 796
monorepo-only commits, \textbf{zero} were present as
\code{(:Commit \{in\_jit:true\})} nodes, and \textbf{zero} had any AST-delta
edge (\code{ADDS}/\code{REMOVES}/\code{UPDATES}/\code{MOVES}). Only the 525
spin-off commits were in the graph at all. The structural layers---the core of
the method---were therefore built on 40\% of the intended stream.

\subsection{Temporal structure of the mix}
Although the two repositories jointly cover overlapping calendar years, the
\emph{labelled} commit stream is almost perfectly segmented in time. Ordering
the 1{,}321 labelled commits by author date and marking each as spin-off (S) or
monorepo (M) yields only \textbf{15 S$\leftrightarrow$M transitions} over the
whole stream:

\begin{center}
\begin{tabular}{@{}llrl@{}}
\toprule
Repo & Index range & $n$ & Date span \\
\midrule
S & [0000:0141]   & 142 & 2009-05-19 \ldots{} 2009-09-07 \\
M & [0142:0142]   & 1   & 2009-09-07 \\
S & [0143:0188]   & 46  & 2009-09-07 \ldots{} 2009-09-19 \\
M & [0189:0189]   & 1   & 2009-09-19 \\
S & [0190:0331]   & 142 & 2009-09-19 \ldots{} 2010-02-22 \\
M & [0332:0332]   & 1   & 2010-02-25 \\
S & [0333:0365]   & 33  & 2010-02-26 \ldots{} 2010-04-21 \\
M & [0366:0366]   & 1   & 2010-04-22 \\
S & [0367:0380]   & 14  & 2010-04-25 \ldots{} 2010-05-12 \\
M & [0381:0381]   & 1   & 2010-05-14 \\
S & [0382:0479]   & 98  & 2010-05-20 \ldots{} 2011-02-05 \\
M & [0480:0480]   & 1   & 2011-02-09 \\
S & [0481:0506]   & 26  & 2011-02-18 \ldots{} 2011-05-10 \\
M & [0507:0507]   & 1   & 2011-05-10 \\
S & [0508:0531]   & 24  & 2011-05-12 \ldots{} 2011-06-12 \\
M & [0532:1320]   & 789 & 2011-06-13 \ldots{} 2019-11-27 \\
\bottomrule
\end{tabular}
\end{center}

The stream is essentially: a long spin-off era (indices 0--531, seven lone
monorepo commits sprinkled in), then a long pure-monorepo era (indices
532--1320). This near-clean segmentation is what makes the fix low-risk (see
\S\ref{sec:fix}).

% ======================================================================
\section{Decision}
% ======================================================================
Three remedies were considered:

\begin{enumerate}
  \item \textbf{Monorepo-only rebuild.} Point the build at \code{apache/hadoop}
  and keep only the 796 commits that resolve there. \emph{Rejected}: it discards
  the 525 spin-off-only commits (the entire 2009--2011 early era, 40\% of the
  labelled set), silently redefining what ``hadoop-mapreduce'' means relative to
  the other ApacheJIT projects and biasing every cross-project comparison.

  \item \textbf{SHA remapping.} Attempt to locate each spin-off commit inside the
  monorepo under its rewritten SHA (via patch-id or message/tree matching) and
  relabel, then build entirely from the monorepo. \emph{Rejected as the primary
  plan}: history rewriting during the re-merge does not guarantee a 1:1
  patch-id mapping, so this risks dropping or mis-associating labels---a subtle,
  hard-to-audit failure precisely in the region we are trying to repair.

  \item \textbf{Multi-repository build (chosen).} Make the build engine resolve
  each commit from whichever repository contains it, and rebuild on the full
  1{,}321-commit stream. This is faithful to ApacheJIT's labelled set exactly as
  published (same SHAs, same commits, none dropped, none relabelled), and it is a
  one-time engine capability that also benefits any future split-history project.
\end{enumerate}

\textbf{Decision: adopt the multi-repository build.} Fidelity to the published
labelled set is the decisive criterion---the experiments must be run on the
commits ApacheJIT actually labelled, not on a convenient subset or a remapped
approximation.

% ======================================================================
\section{The fix and its logic}
\label{sec:fix}
% ======================================================================

\subsection{What the build engine does per commit}
For each labelled commit, in chronological order, the online engine maintains a
\code{file\_state} map: for every tracked \code{.java} file it stores the commit
SHA at which that file was last seen. When a commit modifies a file, the engine
diffs the file's \emph{previous} blob against its \emph{current} blob:

\begin{lstlisting}
before = git_bytes(file_state[path], path)   # file's PREVIOUS commit SHA
after  = git_bytes(commit,          path)    # this commit's SHA
\end{lstlisting}

Both \code{git\_bytes(ref, path)} and \code{changed\_java\_files(commit)} issue
\code{git -C REPO\_PATH <...>} against a \emph{single} repository. This is why a
naive ``pick the right repo per commit'' is insufficient: the \code{before} blob
may live in a \emph{different} repository than the \code{after} blob at the
$\approx$7 spin-off/monorepo boundaries.

\subsection{Multi-repository resolution}
The fix introduces an \emph{ordered list} of repositories for a project and makes
git access repository-aware:

\begin{itemize}
  \item \textbf{Config.} A project may declare additional repositories. For
  \code{hadoop-mapreduce}:
  \begin{lstlisting}
repo_subpath:  repos/apache/hadoop-mapreduce  # primary: base + 525 early commits
extra_repo_subpaths:
  - repos/apache/hadoop                       # monorepo: 796 later commits
  \end{lstlisting}
  \code{project\_config} exposes these as an ordered \code{REPO\_PATHS} list
  (primary first). Single-repository projects are unaffected: with no
  \code{extra\_repo\_subpaths}, \code{REPO\_PATHS == [REPO\_PATH]} and behaviour
  is byte-for-byte identical to before.

  The \emph{spin-off} is chosen as primary because it owns the
  \code{base\_commit} (\code{867a8565}, the earliest labelled commit,
  2009-05-19): Phase~A builds the base AST snapshot by physically
  \code{git checkout <base>} inside the primary repo and reading the working
  tree, so the base must live there. The online growth phase then draws the 796
  later commits from the monorepo via the fallback resolver.

  \item \textbf{Commit-owning repo.} A helper resolves, and caches, which
  repository contains a given commit SHA (\code{git cat-file -e <sha>}), so
  \code{changed\_java\_files(commit)} runs against the repo that owns
  \emph{that} commit.

  \item \textbf{Blob fallback.} \code{git\_bytes(ref, path)} tries each
  repository in \code{REPO\_PATHS} until \code{<ref>:<path>} resolves. In the
  $>98\%$ of the stream that is a contiguous single-repo run, the first repo
  hits immediately; only at the handful of boundaries does the fallback engage.
  If a prior blob genuinely resolves in no repository (a possible boundary
  artifact), the existing code already treats \code{before is None} as a safe
  skip---no crash, no corruption.
\end{itemize}

\subsection{Base ingestion (Phase A2) is multi-repository too}
A subtlety that was \emph{not} obvious at first: making only the online-growth
git access (B5/B6) multi-repository is insufficient. The base ingestion step
(A2, \code{ingest\_base\_kg\_and\_ast.py}) creates the \code{:Commit} nodes by
streaming a repository's \code{git log}, and the labelling step (A4) only sets
\code{in\_jit=true} on commits that \emph{already exist} as nodes
(\code{MATCH (c:Commit \{id:...\})}). If A2 ingests only one repository, the 796
monorepo commits never become nodes, A4 cannot label them, and the online layers
then grow over just the 525 that do exist---exactly the original defect. A first
rebuild attempt reproduced this: the graph still held only 525
\code{in\_jit} commits.

The fix therefore also makes A2 multi-repository, with an asymmetric,
cost-aware policy:
\begin{itemize}
  \item the \textbf{primary} repository (the spin-off) is streamed in full, as
  before, so the base graph keeps its unlabelled context commits;
  \item each \textbf{secondary} repository (the monorepo) is streamed
  \emph{labelled-commits-only}, via \code{git log --no-walk --stdin} fed the
  labelled SHA set (one per line). This avoids ingesting the monorepo's
  $\approx$28{,}000 unrelated commits (the project only needs its 796 labelled
  ones from there). The SHAs are passed on \code{--stdin} rather than as
  command-line arguments because 796 hashes ($\approx$32\,KB) exceed Windows'
  command-line length limit (\code{WinError 206});
  \item commits are de-duplicated by id across repositories (primary wins).
\end{itemize}
A2 now prints an explicit labelled-coverage line
(\code{coverage 1321/1321}) and warns if any labelled commit is missing from all
configured repositories---the guardrail called for in \S8.

\subsection{Why this is low-risk}
The near-clean temporal segmentation (\S3.4, only 15 transitions) means the
stateful \code{before}/\code{after} diffing almost always stays within one
repository. The cross-repo fallback is exercised only around $\approx$7 lone
boundary commits, and its worst case degrades to the pre-existing, already-safe
``skip this file's delta'' path. No change is made to the inference protocol, the
fusion selection, or the scalability methodology; only the \emph{source of git
blobs} during the build becomes multi-repository.

\subsection{Ancillary fix: diff text}
The \code{diff\_text} column of the per-project diff CSV was refilled by resolving
each commit against whichever repository contains it (all 1{,}321 rows now
populated). This restores full input coverage for the CSTG layer. The refill used
the same local-git extraction logic as the original
\code{scripts/build\_diffs.py} (no network calls).

% ======================================================================
\section{Consequences and required re-runs}
% ======================================================================
The \code{hadoop-mapreduce} artifacts produced before this fix are
\textbf{invalid} and must be regenerated:

\begin{itemize}
  \item the Neo4j graph and its snapshot
  (\code{outputs/hadoop-mapreduce/neo4j\_dump/}) were built on 525/1{,}321
  commits;
  \item the experiment results (\code{final\_experiments\_results.pkl},
  \code{final\_fusion\_results.pkl}, subgraph RQ results), the chosen fusion
  (previously reported as DW+CSTG), the scalability suite, and all rendered
  tables/figures/notebooks derive from that partial graph.
\end{itemize}

The full remedy is therefore: rebuild all layers on the complete 1{,}321-commit
multi-repository stream, then re-run experiments, scalability, reporting, and
re-snapshot.

% ======================================================================
\section{Lessons and guardrails}
\label{sec:guardrails}
% ======================================================================
\begin{enumerate}
  \item \textbf{Coverage assertion.} The build should assert that the number of
  \code{(:Commit \{in\_jit:true\})} nodes equals the number of labelled commits,
  and fail loudly on a shortfall, rather than reporting success. A commit whose
  SHA does not resolve in any configured repository must be an error, not a
  silent skip.

  \item \textbf{Provenance check before building.} For any new project, verify up
  front that every labelled SHA resolves in the configured repositories (the
  525/796 partition would have been caught immediately by such a check).

  \item \textbf{Repository-split awareness.} ApacheJIT spans many Apache projects,
  several of which underwent repository splits/merges. Any project whose labelled
  SHAs do not all resolve in a single clone is a candidate for the same
  multi-repository treatment.
\end{enumerate}

% ======================================================================
\section{Second multi-repository case: \code{hadoop-hdfs}}
\label{sec:hdfs}
% ======================================================================
The guardrail from \S\ref{sec:guardrails} above (\emph{provenance check before
building}) paid off immediately on the next Hadoop sub-project. \code{hadoop-hdfs}
is the sibling of \code{hadoop-mapreduce}: it went through the same 2009--2011
spin-off/re-merge, so its ApacheJIT labels are likewise split across the
\code{apache/hadoop-hdfs} spin-off and the \code{apache/hadoop} monorepo. This
time the split was detected \emph{up front} by resolving every labelled SHA
against both clones \emph{before} any build was started---so, unlike
\code{hadoop-mapreduce}, no invalid single-repository graph was ever produced.

\subsection{The label partition}
Resolving each of the 2{,}907 \code{hadoop-hdfs} labelled SHAs against both
repositories:

\begin{center}
\begin{tabular}{@{}lr@{}}
\toprule
Resolves in \ldots{} & Commits \\
\midrule
monorepo \emph{only}            & 2{,}304 \\
spin-off repository \emph{only} & 461 \\
both repositories               & 142 \\
neither repository              & 0 \\
\midrule
total labelled                  & 2{,}907 \\
\bottomrule
\end{tabular}
\end{center}

\subsection{Two differences from \code{hadoop-mapreduce}}
The same multi-repository engine (\S\ref{sec:fix}) handles \code{hadoop-hdfs}
unchanged, but two properties differ and are worth recording:

\begin{itemize}
  \item \textbf{Non-empty intersection.} Here \textbf{142} commits resolve in
  \emph{both} repositories (for \code{hadoop-mapreduce} the intersection was
  empty). This is harmless: \code{git\_root\_for} probes \code{REPO\_PATHS} in
  order and returns the first repository that resolves the SHA, so a
  both-repositories commit is deterministically assigned to the primary. The
  \code{before}/\code{after} blob fallback still works because both clones share
  those objects.

  \item \textbf{The spin-off must be primary --- because primary is streamed in
  full.} Recall (\S\ref{sec:fix}) that A2 streams the \emph{primary} repository's
  \emph{entire} \code{git log} (to keep unlabelled context commits) and streams
  each \emph{secondary} repository labelled-commits-only. The primary must
  therefore be the \emph{small} repository. A first \code{hadoop-hdfs} attempt
  set the \emph{monorepo} primary (it resolves the majority of the stream and the
  base commit), which seemed reasonable --- but full-streaming the
  $\approx$28{,}289-commit monorepo with \code{git log -p --numstat --summary} was
  prohibitively slow (16{,}000 commits in $\sim$76 minutes, and climbing) and
  would have ingested tens of thousands of irrelevant commits. The build was
  stopped and the configuration corrected to make the \emph{spin-off} primary
  (1{,}134 commits, full-stream cheap), with the monorepo as the labelled-only
  extra. The base commit (\code{4ba4f741d159e8\ldots}, 2009-05-19, 277 \code{.java}
  files) resolves in \emph{both}, so Phase~A's \code{git checkout <base>} works on
  the spin-off.
  \begin{lstlisting}
repo_subpath:  repos/apache/hadoop-hdfs  # primary: spin-off (SMALL; full-stream)
extra_repo_subpaths:
  - repos/apache/hadoop                  # monorepo: labelled-only (2,304)
  \end{lstlisting}
  \textbf{Rule of thumb:} the primary repository is whichever \emph{small} clone
  owns (or shares) the base commit; the large monorepo is always a
  \emph{labelled-only} secondary. Choosing primary by ``where most commits
  resolve'' is a trap when that repository is a giant monorepo.
\end{itemize}

\subsection{Ancillary fix and a base-commit transcription trap}
As with \code{hadoop-mapreduce}, the \code{diff\_text} column was refilled by
resolving each commit against whichever repository contains it: 2{,}304
previously-empty rows were populated (all 2{,}907 now complete), restoring CSTG
input coverage.

One process note for future setups: an initial \code{base\_commit} value was
mis-transcribed by hand and resolved in \emph{neither} clone. Because
\code{git\_root\_for(BASE\_COMMIT)} returned \code{None} during a pre-build
configuration check, the error was caught \emph{before} the (multi-hour) Phase~A
build rather than after---a concrete instance of why the provenance check
(\S\ref{sec:guardrails}) belongs \emph{before} the build. Always derive the base SHA
programmatically (\code{find\_base\_commit.py}) or copy it verbatim, and confirm
\code{git\_root\_for(BASE\_COMMIT)} is non-\code{None} before building.

% ======================================================================
\section{A performance wall (not a defect): AST-write cost on a large graph}
\label{sec:perfwall}
% ======================================================================
While building \code{hadoop-hdfs}, the AST-delta layer (B5) slowed to a crawl
around commit 2{,}227 of 2{,}907, with individual commits taking \emph{minutes}
each. This section records the investigation because the symptom (a build that
looks ``stuck'') is alarming but the cause is a well-understood graph-database
performance issue, \emph{not} a correctness defect in the multi-repository work.

\subsection{First question: is the node count itself wrong?}
The alarm was sharpened by an apparent anomaly: at only 76\% built,
\code{hadoop-hdfs} already held \textbf{5.19\,M} \code{:ASTNode}s, \emph{more}
than the fully-built \code{groovy} graph (3.70\,M) despite groovy having
$\sim$3.6$\times$ more commits (8{,}059 vs 2{,}907). A natural worry was that the
multi-repository build was \emph{double-counting}---e.g. re-bootstrapping whole
file ASTs at each spin-off/monorepo boundary.

That worry was tested directly and \textbf{rejected}. A ``bootstrap'' (the engine
capturing an entire file's AST the first time a commit touches a file it is not
yet tracking) is a normal part of the online-growth design, not a boundary
artifact. Tagging each processed commit by its owning repository and correlating
with bootstrap events showed the bootstraps are \emph{not} concentrated at the
repo transitions:

\begin{center}
\begin{tabular}{@{}lr@{}}
\toprule
Bootstrap distribution (2{,}227 commits processed) & Count \\
\midrule
commits with $\geq 1$ bootstrap                 & 622 \\
\quad occurring \emph{at} a repo transition     & 25 \\
\quad occurring \emph{within} a same-repo run   & 597 \\
repo transitions in the stream                  & 55 \\
\bottomrule
\end{tabular}
\end{center}

Bootstraps begin at commit index~3 (a pure spin-off commit, well before any
boundary) and are spread throughout. The high node count is therefore
\emph{genuine}, and is explained by code volume, not a bug:

\begin{center}
\begin{tabular}{@{}lrr@{}}
\toprule
 & \code{groovy} (done) & \code{hadoop-hdfs} (76\%) \\
\midrule
\code{:ASTNode}s            & 3{,}699{,}439 & 5{,}194{,}906 \\
distinct files              & 3{,}012       & 3{,}323 \\
nodes / file                & 1{,}228       & 1{,}563 \\
nodes / commit              & 459           & 2{,}333 \\
\bottomrule
\end{tabular}
\end{center}

Two legitimate factors compound: \code{hadoop-hdfs} files are $\sim$27\% larger
in AST terms (enterprise Java, bigger classes), and---dominantly---its commits
touch $\sim$5$\times$ more AST \emph{per commit} (larger, more file-spanning
changes). \textbf{Graph size tracks code volume changed, not commit count.} Of
the 5.19\,M nodes, 3.42\,M (66\%) are full-file bootstrap nodes and 1.77\,M are
delta nodes---consistent with a large, heavily-churned codebase. This is itself a
useful scalability finding for RQ2.

\subsection{Root cause of the slowdown: an unindexed liveness filter}
The per-commit slowness (a single query observed running for
\code{PT1M38S}) was traced to the two hot queries in
\code{build\_online\_kg.py}---\code{load\_current\_ast} and \code{delete\_file}:

\begin{lstlisting}
MATCH (c:Commit {id:$cid}), (a:ASTNode {file:$file})
WHERE coalesce(a.alive, true)
...
\end{lstlisting}

An index existed on \code{ASTNode(file, pos\_line, pos\_col)}, which serves the
\code{file} equality but returns \emph{all} nodes ever created for that
file---including the many \emph{dead} versions accumulated over thousands of
edits---and then filters \code{alive} in memory. For a hot HDFS file with
hundreds of thousands of historical nodes, that in-memory filter dominates, and
its cost \emph{grows with graph size}: exactly the ``missing-index'' cliff. This
is a known Neo4j behaviour, independent of the multi-repository change (a
single-repo project of the same size would hit it too).

\subsection{The fix: a composite \code{(file, alive)} index}
A composite index lets the planner filter both properties in the index and return
only the live nodes for a file:

\begin{lstlisting}
CREATE INDEX astnode_file_alive IF NOT EXISTS
FOR (a:ASTNode) ON (a.file, a.alive)
\end{lstlisting}

\textbf{Critically, the index alone is not enough---the query must also change.}
The original predicate \code{coalesce(a.alive, true)} \emph{defeats} the index:
Neo4j cannot turn a \code{coalesce(...)} wrapper into an index seek, so the
planner falls back to a label/property scan even with the composite index
present. The predicate must be the bare \code{a.alive = true} for the planner to
choose \code{NodeIndexSeek}. This was confirmed by direct measurement on the
hottest file in the resident graph (141{,}616 accumulated nodes; 42{,}540 live):

\begin{center}
\begin{tabular}{@{}lrl@{}}
\toprule
Predicate & Time & Plan \\
\midrule
\code{coalesce(a.alive, true)} & 67{,}214\,ms & scan (index unused) \\
\code{a.alive = true}          & 437\,ms     & \code{NodeIndexSeek} \\
\bottomrule
\end{tabular}
\end{center}

A \textbf{154$\times$} speedup \emph{on this single isolated query} (the most-edited
HDFS file had accumulated \textbf{141{,}616} AST-node versions over its history,
which is what the old predicate scanned on every touch). This query runs once per
changed file per commit, so it is \emph{a} major cost---but, importantly, not the
only one (see \S\ref{sec:writecost}). Measured on the build itself (first
$\sim$100 commits of the AST layer, like-for-like against the pre-fix run), the
index fix roughly halved end-to-end per-commit time and, more importantly,
\emph{eliminated the catastrophic tail}:

\begin{center}
\begin{tabular}{@{}lrrr@{}}
\toprule
AST per-commit wall & median & mean & worst \\
\midrule
before any fix & $4.0$\,s & $12.8$\,s & $300$--$1130$\,s \\
with index $+$ \code{a.alive=true} & $2.75$\,s & $5.5$\,s & $\sim$$60$\,s \\
\bottomrule
\end{tabular}
\end{center}

The tail collapse (worst commit $\sim$19\,min $\to$ $\sim$1\,min) is the decisive
practical win: it turns an effectively-unbounded build into a finite one.

The complete fix has three parts, all now in the build engines:
\begin{itemize}
  \item \textbf{Rewrite the predicate} from \code{coalesce(a.alive, true)} to
  \code{a.alive = true} in every hot per-file liveness query
  (\code{load\_current\_ast}, \code{delete\_file} in both
  \code{build\_online\_kg.py} and \code{build\_subgraph\_online\_kg.py}; the
  one-off report/count queries in \code{collect\_subgraph\_stats.py},
  \code{ingest\_cstg.py}, \code{collect\_kg\_profile.py} were simplified for
  consistency).
  \item \textbf{Backfill} any \code{alive IS NULL} nodes to \code{true} at build
  startup (base nodes from A2 may lack the property; there were 859 here). With
  no nulls, \code{a.alive = true} is exactly equivalent to the old
  \code{coalesce} semantics but index-seekable.
  \item \textbf{Create the composite index at startup} via
  \code{CREATE INDEX IF NOT EXISTS}, so every large project---\code{camel}
  (22.7\,k commits), \code{ignite}, \code{flink}---benefits automatically rather
  than hitting the same wall. The subgraph engine creates the per-label
  equivalent (\code{<label>\_file\_alive}).
\end{itemize}

\subsection{A write-path optimisation that was attempted and reverted}
\label{sec:writecost}
Beyond the read query, a second cost was observed: a per-commit write transaction
in \code{[Closing]} state holding $\approx$$50{,}974$ locks for $\sim$2 minutes.
The AST inserts use \code{MERGE (a:ASTNode \{id:\ldots\})}, and since an
\code{ASTNode.id} \emph{uniqueness constraint} already exists, the MERGE's
read-before-write $+$ lock looked redundant; \code{CREATE} promised to be much
cheaper. This was implemented (with per-commit checkpointing to keep resumes
safe) \textbf{and then reverted}, for an instructive reason:

\begin{quote}
The base-style AST ids \code{\{file\}::A\{i\}} are deterministic \emph{per file},
not per commit. A file that is deleted and re-added, or bootstrapped via two code
paths, legitimately re-touches the same ids. \code{MERGE} is idempotent there;
\code{CREATE} raised \code{Neo.ClientError.Schema.ConstraintValidationFailed} on a
clean build. Only the \emph{delta} inserts (\code{\{file\}::D\{sha\}:\{i\}},
commit-stamped) have truly-unique ids---too narrow a slice to justify the
correctness risk. The engines were restored to \code{MERGE} everywhere; only the
read-query fix above was kept.
\end{quote}

The lesson: an ``obvious'' \code{MERGE}$\to$\code{CREATE} swap is unsafe whenever
node ids can legitimately recur; verify id lifetime before assuming uniqueness.
The residual write cost remains, but with the read fix the build is finite and the
tail is bounded, which was the blocking problem.

\subsection{Operational lessons}
\begin{itemize}
  \item \textbf{Never bulk-\code{DELETE} a multi-million-node layer in one
  transaction.} Running \code{build\_online\_kg.py --reset} on the resident
  $5.2$\,M-node graph issued a single Cypher delete that \textbf{OOM-crashed
  Neo4j} (\code{java.lang.OutOfMemoryError}, Bolt refusing connections). The graph
  survived (the transaction rolled back on restart), but the reset did not
  complete. To wipe a large project's store, use the \emph{store-level} wipe
  (\code{reset\_neo4j.py}: stop container, \code{rm -rf} the DB files, restart),
  which is $O(1)$ in memory, not a transaction.
  \item \textbf{Transient Bolt drops on heavy commits are recoverable.} A drop
  mid-write (\code{\_buffer\_one\_chunk}) leaves the checkpoint intact; the layer
  resumes from \code{next\_index}. With the reverted \code{MERGE} engine, resume is
  idempotent, so re-processing the last checkpoint window is safe; do \emph{not}
  pass \code{--reset} to a resume.
  \item \textbf{Measure the build, not an isolated query, and clear stale logs
  first.} An early ``speedup'' reading was actually a \emph{stale}
  \code{ast\_timing.csv} from a prior run; the real like-for-like comparison
  (above) required deleting the log before re-measuring. Isolated-query speedups
  ($154\times$) do not translate one-to-one to end-to-end build time.
  \item \textbf{``Stuck'' vs ``slow'' is decided by the checkpoint, not the log.}
  Forward progress is proven by \code{next\_index} advancing and by inspecting
  \code{SHOW TRANSACTIONS} (elapsed time, \code{activeLockCount}) for a
  long-running query---used here to tell a genuinely slow write from a hang.
\end{itemize}

\end{document}
